\documentclass[12pt]{article}
\usepackage[utf8]{inputenc}
\usepackage[T1]{fontenc}
\usepackage{hyperref}
\title{Dorig Language Resources for K-12 Curriculum Development}
\author{Compiled by Claude Sonnet 5.5 (claude-sonnet-5-5) for Andr\'as Kornai}
\date{}
\begin{document}
\maketitle

\section{Introduction}
Dorig (ISO 639-3: wwo; formerly called Wetamut) is an Oceanic (Austronesian) language of Vanuatu, spoken on the south coast of Gaua island in the Banks Islands. It takes its name from the village of D\=orig. Smaller groups of speakers live in Qteon (east coast) and Qtevut (west coast). Estimates are about 300 speakers (2012). UNESCO classifies it as definitely endangered. Searches found descriptive and catalogue material only, with no school materials.

\section{Primary Educational Resources}
No primers, textbooks or school programmes for Dorig turned up in searches. The sources note that the language has no literacy tradition.

\section{Linguistic and Reference Materials}
\begin{itemize}
\item Wikipedia article, covering the phonology (eight vowels, 15 consonants) and the other languages of Gaua: \url{https://en.wikipedia.org/wiki/Dorig_language}.
\item Endangered Languages Project entry, with a bibliography: \url{https://endangeredlanguages.com/lang/2427}.
\item Omniglot page: \url{https://www.omniglot.com/writing/dorig.htm}.
\end{itemize}

\section{Religious and Community Resources}
No scripture or community-organisation resources were found in the search results.

\section{Digital and Audio Resources}
The Pangloss Collection (CNRS) page on the languages of Gaua island may include recordings, but I did not check its contents: \url{https://pangloss.cnrs.fr/corpus/Nume?lang=en}. It is listed as a corpus lead only.

\section{Curriculum Development Context}
Dorig is one of six closely related languages on Gaua (Nume, Mwerlap, Dorig, Koro, Olrat and Lakon), so materials could be developed jointly across the island. The language has no literacy tradition and little published description. A first step would be community-led work on an orthography, a dictionary and recorded stories, drawing on the linguistic literature listed in the Endangered Languages Project bibliography. Vanuatu schools use English and French, with Bislama widely spoken, so Dorig would most likely start as a community or early-primary subject.

\end{document}